\documentclass[12pt]{letter}

\usepackage{amsmath}
\usepackage{amssymb}

\usepackage{nopageno}
\usepackage{amsmath}
\usepackage{enumitem}
% \usepackage{mathtools}
\usepackage{eqnarray,amsmath}
\usepackage[a4paper, total={7.5in, 11.25in}]{geometry}

\begin{document}

Name: \hrulefill \quad Neptun code: \hrulefill

\begin{center}
(KTXFI2EBNF) Physics II. Test 2 (Midterm Retake) \\
December 8, 2025
\end{center}

\bigskip

\begin{enumerate}
\setcounter{enumi}{0}

\item An electron has a kinetic energy of $0.5\,\text{keV}$. \\

\begin{enumerate}
\item What is its velocity?\\

Convert to SI: $K = 0.5\,\text{keV} = 500\,\text{eV} = 500\cdot(1.602176634\cdot10^{-19})\,\text{J} \approx 8.01088317\cdot10^{-17}\,\text{J}.$
Using the nonrelativistic kinetic energy formula
$$ K=\frac{1}{2} m v^{2}, $$
with $m=m_e=9.1\cdot10^{-31}\,\text{kg}$ gives
$$ v=\sqrt{\frac{2K}{m}}\approx \sqrt{\frac{2\cdot 8.01088317\cdot10^{-17}}{9.1\cdot10^{-31}}}
\approx 1.3269\cdot10^{7}\,\text{m/s}. $$

\item What is its momentum?

$$ p=m v \approx (9.1\cdot10^{-31}\,\text{kg})\cdot(1.3269\cdot10^{7}\,\text{m/s})
\approx 1.2075\cdot10^{-23}\,\text{kg\,m/s}. $$

\item What is the wavelength of the associated de Broglie wave?\\

Using $\lambda = h/p$ with $h=6.62607015\cdot10^{-34}\,\text{J,s}$,
$$ \lambda=\frac{h}{p}\approx \frac{6.62607015\cdot10^{-34}}{1.2075\cdot10^{-23}}
\approx 5.4876\cdot10^{-11}\,\text{m} \approx 0.0549\,\text{nm}. $$
\hfill (10 points)
\end{enumerate}

\bigskip
\item A metal surface is illuminated with light of wavelength $9\cdot10^{-8}\,\text{m}$. What is the velocity of the emitted
electrons if the photoelectric effect starts at a wavelength of $2.88\cdot10^{-7}\,\text{m}$? \
(Use $m_e=9.1\cdot10^{-31}\,\text{kg}$.)\

Photon energy: $E_\text{photon}=hc/\lambda$. The electron kinetic energy equals the excess energy:
$$ K = hc\left(\frac{1}{\lambda_\text{in}}-\frac{1}{\lambda_\text{th}}\right), $$
with $h=6.62607015\cdot10^{-34}\,\text{J,s}$ and $c=2.99792458\cdot10^{8}\,\text{m/s}$. Substituting $\lambda_\text{in}=9.0\cdot10^{-8}\,\text{m}$ and $\lambda_\text{th}=2.88\cdot10^{-7}\,\text{m}$ gives
$$ K \approx 1.5174\cdot10^{-18}\,\text{J}. $$
Then the emitted-electron velocity (nonrelativistic approximation)
$$ v=\sqrt{\frac{2K}{m_e}}\approx \sqrt{\frac{2\cdot1.5174\cdot10^{-18}}{9.1\cdot10^{-31}}}
\approx 1.8262\cdot10^{6}\,\text{m/s}. $$
\hfill (10 points)

\bigskip
\item In a Compton scattering experiment, X-rays with a wavelength of $0.2\,\text{nm}$ are used.\\
($m_e=9.1\cdot10^{-31}\,\text{kg}$)
\begin{enumerate}
\item At what scattering angle does the wavelength of the radiation increase by 1\%?\\

Initial wavelength $\lambda_0=0.2\,\text{nm}=2.00\cdot10^{-10}\,\text{m}$. A 1\% increase means  $\Delta\lambda=0.01\lambda_0=2.00\cdot10^{-12}\,\text{m}.$
The Compton wavelength is
$$ \lambda_C=\frac{h}{m_e c}\approx \frac{6.62607015\cdot10^{-34}}{(9.1\cdot10^{-31})(2.99792458\cdot10^{8})}
\approx 2.4288\cdot10^{-12}\,\text{m}. $$
Compton formula:
$$ \Delta\lambda=\lambda_C(1-\cos\theta). $$
Hence
$$ 1-\cos\theta=\frac{\Delta\lambda}{\lambda_C}\approx\frac{2.00\cdot10^{-12}}{2.4288\cdot10^{-12}}\approx 0.82345, $$
so
$$ \cos\theta \approx 0.17655,\qquad \theta\approx \arccos(0.17655)\approx 79.83^{\circ}. $$

\item At what angle does the wavelength become 0.01% larger?\

Now $\Delta\lambda=0.0001\lambda_0=2.00\cdot10^{-14}\,\text{m}$, therefore
$$ 1-\cos\theta=\frac{2.00\cdot10^{-14}}{2.4288\cdot10^{-12}}\approx 0.0082378, $$
so
$$ \cos\theta \approx 0.9917622,\qquad \theta\approx \arccos(0.9917622)\approx 7.36^\circ. $$
\hfill (10 points)
\end{enumerate}

\bigskip
\item A light source has a power output of $20\,\text{W}$.

\begin{enumerate}
\item What is the frequency of the light if the entire energy of the source is emitted as light with
wavelength $\lambda = 800\,\text{nm}$?\

$$ f=\frac{c}{\lambda}=\frac{2.99792458\cdot10^{8}}{800\cdot10^{-9}}\,\text{Hz}\approx 3.7474\cdot10^{14}\,\text{Hz}. $$

\item  How many photons are emitted by the light source per second?\

Each photon has energy $E=hf$, therefore the photon emission rate is
$$ \dot{N}=\frac{P}{hf}=\frac{P\lambda}{hc}.$$
With $P=20\,\text{W}$ and $\lambda=800\cdot10^{-9}\,\text{m}$,
$$ \dot{N}\approx \frac{20\cdot 800\cdot10^{-9}}{(6.62607015\cdot10^{-34})(2.99792458\cdot10^{8})}
\approx 8.0546\cdot10^{19}\,\text{photons/s}. $$
\hfill (10 points)
\end{enumerate}

\bigskip
\item We study an electron in a radioactive particle beam as it passes through the laboratory. It is found that, as measured in the laboratory, the average lifetime of the particle is $60\,\text{ns}$. The average lifetime measured in the particle rest frame is $5\,\text{ns}$. What is the propagation velocity of the particle beam? Determine the value of $\dfrac{m}{m_0}$ for the electron! Calculate the rest energy of an electron corresponding to its rest mass ($m_e=9.1\cdot10^{-31}\,\text{kg}$) in joules and in electronvolts. \

Time dilation gives $\tau=\gamma\tau_0$, hence
$$ \gamma=\frac{\tau}{\tau_0}=\frac{60\,\text{ns}}{5\,\text{ns}}=12. $$
Therefore
$$ \beta=\frac{v}{c}=\sqrt{1-\frac{1}{\gamma^{2}}}=\sqrt{1-\frac{1}{144}}\approx 0.99652, $$
and
$$ v\approx 0.99652,c \approx 2.9875\cdot10^{8}\,\text{m/s}. $$
The relativistic mass factor is
$$\frac{m}{m_0}=\gamma\approx 12. $$
The rest energy of the electron is
$$ E_{0}=m_e c^{2}\approx (9.1\cdot10^{-31}\,\text{kg})\cdot(2.99792458\cdot10^{8}\,\text{m/s})^{2}
\approx 8.1787\cdot10^{-14}\,\text{J}. $$
In electronvolts,
$$ E_{0} \approx \frac{8.1787\cdot10^{-14}\,\text{J}}{1.602176634\cdot10^{-19}\,\text{J/eV}}
\approx 5.1047\cdot10^{5}\,\text{eV}\approx 510.5\,\text{keV}\approx 0.5105\,\text{MeV}. $$
\hfill (10 points)

\end{enumerate}

% ------------------------------------------------------------------------------------------------------------------------------------

Requirements for obtaining a signature: achieving at least 40% (20,points) on the midterm exam.

Grades: 0–25 points: Fail (1), 26–30 points: Pass (2), 31–37 points: Satisfactory (3), 38–45 points: Good (4), 46–50 points: Excellent (5).

\rightline{Dr.~Gabor FACSKO}
\rightline{\textit{facsko.gabor@uni-obuda.hu}}

\end{document}
